\documentclass[unicode, 12pt, a4paper, oneside, fleqn]{article}

\usepackage{styles/log-style}

\setcounter{secnumdepth}{2}
\setcounter{tocdepth}{2}
\setlength{\LTleft}{0pt}
\setlength{\LTright}{0pt}

\newcommand{\reporttype}{Курсовая работа}
\newcommand{\reportcourse}{Операционные системы}
\newcommand{\reportauthor}{К.\,Р. Шестаков}
\newcommand{\reportgroup}{М8О-207БВ-24}
\newcommand{\reportteacher}{Е.\,С. Миронов}
\newcommand{\reportdate}{22.04.2026}
\newcommand{\reportcity}{Москва}
\newcommand{\reportyear}{2026}
\newcommand{\powertwoname}{2\textasciicircum{}n}

\begin{document}

\hypersetup{pageanchor=false}
\begin{titlepage}
\begin{center}
\large
МОСКОВСКИЙ АВИАЦИОННЫЙ ИНСТИТУТ\\
(НАЦИОНАЛЬНЫЙ ИССЛЕДОВАТЕЛЬСКИЙ УНИВЕРСИТЕТ)

\vspace{20pt}

Институт №8 «Компьютерные науки и прикладная математика»

Кафедра 806 «Вычислительная математика и программирование»
\end{center}

\vspace{58pt}

\begin{center}
{\Large\bfseries \reporttype\par}
\vspace{6pt}
{\large по курсу \reportcourse\par}

\vspace{54pt}
\end{center}

\vfill

\begin{flushright}
\large
\begin{tabular}{rl}
Выполнил: & \reportauthor \\
Группа: & \reportgroup \\
Преподаватель: & \reportteacher \\
Дата: & \reportdate \\
\end{tabular}
\end{flushright}

\vspace{84pt}

\begin{center}
\large
\reportcity, \reportyear
\end{center}
\end{titlepage}

\clearpage
\hypersetup{pageanchor=true}
\pagestyle{plain}
\setcounter{page}{1}

\section{Условие}

\noindent\textbf{Цель работы:}
\begin{enumerate}[leftmargin=1.5cm]
\item реализовать два пользовательских аллокатора памяти, работающих поверх заранее выделенного буфера;
\item сравнить их по скорости выделения памяти, скорости освобождения и эффективности использования памяти на одинаковых сценариях нагрузки.
\end{enumerate}

\noindent\textbf{Задание:} спроектировать общий интерфейс распределителей памяти, реализовать аллокатор на основе списка свободных блоков со стратегией \texttt{first fit} и аллокатор с блоками размера $2^n$, подготовить единый тестовый стенд, выполнить серию измерений и сделать выводы по полученным результатам.

\noindent\textbf{Вариант:} реализовать и сравнить пользовательский аллокатор со списком свободных блоков и аллокатор с блоками размера $2^n$. Исследование проводится на собственных реализациях распределителей памяти без обращения к операционной системе во время работы: оба варианта используют заранее выделенный непрерывный буфер.

\section{Метод решения}

\subsection{Общий подход}

Оба аллокатора работают не с системным менеджером памяти процесса напрямую, а внутри заранее подготовленного массива байтов. Такой подход изолирует исследуемую логику от внешних факторов и позволяет сравнивать именно алгоритмы распределения памяти, а не поведение стандартного аллокатора операционной системы.

Снаружи обе реализации скрыты за общей структурой \texttt{Allocator}. В ней хранится тип аллокатора и указатель на внутреннее состояние. Благодаря этому тестовый стенд использует единый интерфейс и может запускать идентичные сценарии для обоих вариантов без изменения управляющего кода.

\subsection{Аллокатор на основе списка свободных блоков}

В первом варианте память представляется как двусвязный список блоков. Каждый блок содержит служебный заголовок с размером полезной области, флагом занятости и ссылками на соседние блоки. После инициализации весь буфер считается одним большим свободным блоком.

При вызове \texttt{alloc} аллокатор последовательно проходит по списку слева направо и выбирает первый свободный блок достаточного размера. Если выбранный блок заметно больше требуемого размера, он разделяется на две части: первая передаётся пользователю, вторая остаётся свободной. При освобождении памяти блок помечается как свободный и затем при возможности объединяется с соседними свободными блоками.

Преимущество такого подхода состоит в более экономном использовании памяти, поскольку размер выделенного блока близок к запрошенному. Основной недостаток --- зависимость времени поиска от длины и текущей фрагментированности списка свободных блоков.

\subsection{Аллокатор с блоками размера \powertwoname}

Второй вариант построен по идее \enquote{buddy allocator}. Вся память разбивается на блоки, размеры которых равны степеням двойки. Для каждого порядка поддерживается собственный список свободных блоков.

Если требуется блок размера \texttt{S}, аллокатор сначала учитывает размер служебного заголовка, затем выбирает минимальный блок $2^n$, в который запрос помещается. Если свободного блока нужного порядка нет, берётся более крупный блок и делится пополам до тех пор, пока не будет получен блок подходящего размера.

При освобождении памяти вычисляется \enquote{парный} блок того же порядка. Если он тоже свободен, оба блока объединяются в более крупный. Такое слияние продолжается, пока это возможно. Достоинство этого подхода --- стабильное и быстрое время выделения. Недостаток --- внутренняя фрагментация, возникающая из-за округления запроса вверх до степени двойки.

\subsection{Критерии сравнения}

В работе сравниваются четыре показателя:
\begin{enumerate}[leftmargin=1.5cm]
\item фактор использования памяти;
\item скорость выделения памяти;
\item скорость освобождения памяти;
\item простота использования и сопровождения реализации.
\end{enumerate}

Под фактором использования памяти понимается отношение полезно запрошенной памяти к объёму, который аллокатор реально занял внутри буфера:
\[
U = \frac{\text{requested\_bytes}}{\text{reserved\_bytes}}.
\]

Чем ближе значение $U$ к единице, тем меньше потери на служебные заголовки и внутреннюю фрагментацию.

\section{Описание программы}

\subsection{Структура проекта}

Проект разделён на несколько файлов:
\begin{itemize}[leftmargin=1.5cm]
\item \texttt{src/allocator.h} --- общий интерфейс распределителей памяти и структура статистики;
\item \texttt{src/allocator.cpp} --- реализация двух аллокаторов и вспомогательных функций;
\item \texttt{src/benchmark.cpp} --- генерация сценариев нагрузки, запуск измерений и вывод результатов;
\item \texttt{CMakeLists.txt} --- конфигурация сборки проекта;
\item \texttt{report/report.md} --- исходный текст отчёта в Markdown;
\item \texttt{report/report.tex} --- оформленная версия отчёта в формате LaTeX.
\end{itemize}

\subsection{Внешний интерфейс}

Обе реализации используют единый публичный интерфейс:

\begin{lstlisting}[language=C++,caption={Публичный интерфейс аллокаторов},captionpos=b]
Allocator* createFirstFitAllocator(void* realMemory, size_t memorySize);
Allocator* createPowerOfTwoAllocator(void* realMemory, size_t memorySize);
void destroyAllocator(Allocator* allocator);

void* alloc(Allocator* allocator, size_t blockSize);
void freeBlock(Allocator* allocator, void* block);

AllocatorStats getAllocatorStats(const Allocator* allocator);
size_t getAllocatedBlockFootprint(const Allocator* allocator, const void* block);
const char* getAllocatorName(const Allocator* allocator);
\end{lstlisting}

Функция освобождения названа \texttt{freeBlock}, чтобы не пересекаться со стандартной функцией \texttt{free}. Дополнительная функция \texttt{getAllocatedBlockFootprint} используется только в тестовом стенде и позволяет корректно вычислять фактический размер блока, занятого внутри буфера.

\subsection{Ключевые структуры данных}

В реализации используются две разные внутренние схемы хранения состояния:
\begin{itemize}[leftmargin=1.5cm]
\item \texttt{FirstFitBlockHeader} и \texttt{FirstFitState} для варианта со списком свободных блоков;
\item \texttt{BuddyBlockHeader} и \texttt{BuddyState} для варианта с блоками размера \powertwoname.
\end{itemize}

В обоих случаях служебные заголовки размещаются непосредственно внутри рабочего буфера. Это позволяет измерять реальные накладные расходы на хранение метаданных и учитывать их в статистике.

\subsection{Тестовый стенд}

Файл \texttt{src/benchmark.cpp} описывает три сценария нагрузки:
\begin{enumerate}[leftmargin=1.5cm]
\item \textbf{Мелкие блоки.} Запросы от \texttt{16} до \texttt{128} байт.
\item \textbf{Смешанная нагрузка.} Запросы от \texttt{32} до \texttt{4096} байт.
\item \textbf{Фрагментация.} Чередование маленьких и больших блоков с частичным освобождением и последующей случайной нагрузкой.
\end{enumerate}

Для измерения времени используются стандартные средства библиотеки C++: \texttt{std::chrono::steady\_clock}. Повторяемость сценариев обеспечивается фиксированными значениями генератора случайных чисел.

\subsection{Используемые системные средства}

Прямые системные вызовы в программе не используются. Это является частью замысла работы: исследуемые аллокаторы во время тестов не обращаются к операционной системе за дополнительной памятью и не зависят от поведения внешнего распределителя. Сборка выполняется через CMake и MSBuild, а выполнение --- в обычном консольном процессе Windows.

\subsection{Сборка и запуск}

Сборка проекта выполняется командами:

\begin{lstlisting}
cmake -S . -B build
cmake --build build --config Release
\end{lstlisting}

Запуск тестового стенда:

\begin{lstlisting}
build\Release\memory_allocators.exe
\end{lstlisting}

\section{Результаты работы}

\subsection{Сценарий тестирования}

Измерения проводились для трёх сценариев нагрузки, описанных выше. Во всех случаях использовался буфер объёмом \texttt{1 048 576} байт. Каждый сценарий повторялся \texttt{5} раз, после чего вычислялись средние значения времени \texttt{alloc} и \texttt{freeBlock}, среднего фактора использования памяти и его пикового значения.

Сборка выполнялась в режиме \texttt{Release} компилятором \texttt{MSVC 2022}. Такой режим выбран для того, чтобы сравнение отражало поведение реализаций без отладочных накладных расходов.

\subsection{Фрагмент фактического вывода программы}

Ниже приведён лог одного из запусков программы, полученный после сборки текущей версии проекта:

\begin{small}
\begin{alltt}
Сравнение аллокаторов памяти
Размер тестового буфера: 1048576 байт
Количество прогонов каждого сценария: 5

=== Мелкие блоки ===
Случайные блоки размером от 16 до 128 байт, умеренное число
одновременно живых объектов.
Список свободных блоков:
  успешных выделений: 25253
  отказов: 0
  alloc, нс: 546.75
  free, нс: 19.30
  средн. использование: 0.63
  пик: 0.79
Блоки 2\^{}n:
  успешных выделений: 25253
  отказов: 0
  alloc, нс: 24.82
  free, нс: 21.91
  средн. использование: 0.47
  пик: 0.73

=== Смешанная нагрузка ===
Блоки размером от 32 до 4096 байт, нагрузка похожа на реальную
программу общего назначения.
Список свободных блоков:
  успешных выделений: 25126
  отказов: 0
  alloc, нс: 368.16
  free, нс: 25.81
  средн. использование: 0.98
  пик: 0.99
Блоки 2\^{}n:
  успешных выделений: 25126
  отказов: 0
  alloc, нс: 29.67
  free, нс: 24.64
  средн. использование: 0.73
  пик: 0.84

=== Фрагментация ===
Чередование маленьких и больших блоков с частичным освобождением
для создания разрывов в памяти.
Список свободных блоков:
  успешных выделений: 6351
  отказов: 5247
  alloc, нс: 1833.77
  free, нс: 29.41
  средн. использование: 0.97
  пик: 0.98
Блоки 2\^{}n:
  успешных выделений: 7942
  отказов: 3656
  alloc, нс: 28.81
  free, нс: 26.60
  средн. использование: 0.54
  пик: 0.59
\end{alltt}
\end{small}

\subsection{Сводные результаты}

\textbf{Сценарий 1. Мелкие блоки}

\begin{center}
{\scriptsize
\setlength{\tabcolsep}{3pt}
\renewcommand{\arraystretch}{1.15}
\begin{tabular}{|p{3.8cm}|r|r|r|r|r|r|}
\hline
\makecell[l]{Аллокатор} &
\makecell[r]{Усп.\\выд.} &
\makecell[r]{Отказы} &
\makecell[r]{Alloc,\\нс} &
\makecell[r]{Free,\\нс} &
\makecell[r]{Средн.\\$U$} &
\makecell[r]{Пик.\\$U$} \\
\hline
Список свободных блоков & 25253 & 0 & 546.75 & 19.30 & 0.63 & 0.79 \\
\hline
Блоки размера $2^n$ & 25253 & 0 & 24.82 & 21.91 & 0.47 & 0.73 \\
\hline
\end{tabular}
}
\end{center}

На мелких блоках аллокатор с блоками размера $2^n$ заметно быстрее выполняет выделение памяти. Разница почти на порядок определяется тем, что выбор подходящего блока идёт по готовым спискам фиксированных порядков. При этом вариант \texttt{first fit} лучше использует память, так как не требует агрессивного округления размера запроса вверх.

\textbf{Сценарий 2. Смешанная нагрузка}

\begin{center}
{\scriptsize
\setlength{\tabcolsep}{3pt}
\renewcommand{\arraystretch}{1.15}
\begin{tabular}{|p{3.8cm}|r|r|r|r|r|r|}
\hline
\makecell[l]{Аллокатор} &
\makecell[r]{Усп.\\выд.} &
\makecell[r]{Отказы} &
\makecell[r]{Alloc,\\нс} &
\makecell[r]{Free,\\нс} &
\makecell[r]{Средн.\\$U$} &
\makecell[r]{Пик.\\$U$} \\
\hline
Список свободных блоков & 25126 & 0 & 368.16 & 25.81 & 0.98 & 0.99 \\
\hline
Блоки размера $2^n$ & 25126 & 0 & 29.67 & 24.64 & 0.73 & 0.84 \\
\hline
\end{tabular}
}
\end{center}

На смешанной нагрузке соотношение сохраняется. Аллокатор с блоками размера $2^n$ остаётся быстрее, однако теряет в эффективности использования памяти. Для \texttt{first fit} особенно хорошо видно, что почти весь фактически занятый объём идёт в полезные данные.

\textbf{Сценарий 3. Фрагментация}

\begin{center}
{\scriptsize
\setlength{\tabcolsep}{3pt}
\renewcommand{\arraystretch}{1.15}
\begin{tabular}{|p{3.8cm}|r|r|r|r|r|r|}
\hline
\makecell[l]{Аллокатор} &
\makecell[r]{Усп.\\выд.} &
\makecell[r]{Отказы} &
\makecell[r]{Alloc,\\нс} &
\makecell[r]{Free,\\нс} &
\makecell[r]{Средн.\\$U$} &
\makecell[r]{Пик.\\$U$} \\
\hline
Список свободных блоков & 6351 & 5247 & 1833.77 & 29.41 & 0.97 & 0.98 \\
\hline
Блоки размера $2^n$ & 7942 & 3656 & 28.81 & 26.60 & 0.54 & 0.59 \\
\hline
\end{tabular}
}
\end{center}

При фрагментации различие между подходами становится ещё сильнее. Аллокатор \texttt{first fit} по-прежнему экономно расходует память, но резко проигрывает по времени выделения, поскольку поиск первого подходящего блока выполняется на длинной и фрагментированной цепочке. Вариант с блоками размера $2^n$ сохраняет почти постоянное время \texttt{alloc} и показывает меньше отказов.

\subsection{Итог сравнения}

По результатам измерений можно сделать следующие выводы:
\begin{itemize}[leftmargin=1.5cm]
\item аллокатор \texttt{first fit} выигрывает по эффективности использования памяти;
\item аллокатор с блоками размера $2^n$ выигрывает по скорости выделения памяти;
\item по скорости освобождения оба подхода близки, но вариант $2^n$ остаётся немного стабильнее на сложных сценариях;
\item с точки зрения объяснимости и сопровождения реализация списка свободных блоков проще, чем реализация \enquote{buddy allocator}.
\end{itemize}

\section{Выводы}

В ходе работы были реализованы два пользовательских аллокатора памяти поверх заранее выделенного буфера и проведено их сравнение на нескольких сценариях нагрузки. Подход на основе списка свободных блоков со стратегией \texttt{first fit} показал более экономный расход памяти и высокую эффективность при умеренной фрагментации.

Аллокатор с блоками размера $2^n$ продемонстрировал существенно более высокую скорость выделения памяти и лучшую устойчивость к фрагментации, однако платой за это стала внутренняя фрагментация и менее высокий коэффициент использования памяти.

Следовательно, однозначно лучшего варианта здесь нет. Если в прикладной задаче важнее минимизировать потери памяти, предпочтителен \texttt{first fit}. Если критична скорость выделения и предсказуемость времени работы даже в неблагоприятных состояниях памяти, практичнее использовать аллокатор с блоками размера $2^n$.

\clearpage
\section{Исходный код программы}

\subsection{Файл \texttt{src/allocator.h}}

\begin{center}
\lstinputlisting[language=C++,caption={src/allocator.h},captionpos=b]{listings/allocator.h.txt}
\end{center}

\subsection{Файл \texttt{src/allocator.cpp}}

\begin{center}
\lstinputlisting[language=C++,caption={src/allocator.cpp},captionpos=b]{listings/allocator.cpp.txt}
\end{center}

\subsection{Файл \texttt{src/benchmark.cpp}}

\begin{center}
\lstinputlisting[language=C++,caption={src/benchmark.cpp},captionpos=b]{listings/benchmark.cpp.txt}
\end{center}

\subsection{Файл \texttt{CMakeLists.txt}}

\begin{center}
\lstinputlisting[language={},caption={CMakeLists.txt},captionpos=b]{listings/CMakeLists.txt}
\end{center}

\end{document}
